其中，例如，$P[Z\leqslant0.385]=\int_{-\infty}^{0.385}\frac{1}{\sqrt{2\pi}}e^{-z^{2}/2}dz=0.65$ [参见\hyperref[eq:4-30]{式(4-30)}].

让我们构造 $Q$-$Q$ 图并评述其形状. 对于上述数据，该 $Q$-$Q$ 图描述的是顺序数据 $x_{i(j)}$ 与正态分位数 $q_{i(j)}$ 之间的关系，如\hyperref[fig:4.5]{图 4.5} 所示. 点对 $(q_{i(j)},x_{i(j)})$ 的位置很接近一条直线，并且我们不拒绝这些数据服从正态分布的看法——特别是在样本容量像 $n=10$ 这么小的时候.

\begin{mvfigure}
\begin{tikzpicture}[scale=0.85]
\draw[->] (-2.9,0) -- (3.0,0) node[below right] {$q_{i(j)}$};
\draw[->] (0,-1.7) -- (0,3.0) node[left] {$x_{i(j)}$};
\foreach \x in {-2,-1,1,2} {\draw (\x,0.08) -- (\x,-0.08) node[below] {\x};}
\node[below left] at (0,0) {$0$};
\foreach \y in {-1,1,2} {\draw (0.08,\y) -- (-0.08,\y) node[left] {\y};}
\draw (-2.0,-1.30) -- (2.0,2.72);
\foreach \p in {(-1.645,-1.00),(-1.036,-0.10),(-0.674,0.16),(-0.385,0.41),(-0.125,0.62),(0.125,0.80),(0.385,1.26),(0.674,1.54),(1.036,1.71),(1.645,2.30)}
  {\fill \p circle (2.2pt);}
\end{tikzpicture}
\phantomsection\label{fig:4.5}\mvcap{图 4.5\quad 例 4.9 中数据的 $Q$-$Q$ 图 \hfill$\blacksquare$}
\end{mvfigure}

为了画 $Q$-$Q$ 图所需要的计算，用计算机很容易编程. 在市场上有许多统计程序能画出这种图形.

作 $Q$-$Q$ 图的步骤是：

1. 把原始观测值排序，得到 $x_{i(1)},x_{i(2)},\cdots,x_{i(n)}$ 和它们的对应概率值 $\left(1-\dfrac{1}{2}\right)\big/n,\left(2-\dfrac{1}{2}\right)\big/n,\cdots,\left(n-\dfrac{1}{2}\right)\big/n$;

2. 计算标准正态分位数 $q_{i(1)},q_{i(2)},\cdots,q_{i(n)}$;

3. 作观测值的点对 $(q_{i(1)},x_{i(1)}),(q_{i(2)},x_{i(2)}),\cdots,(q_{i(n)},x_{i(n)})$ 的图形，并检查结果的“直线性”.

$Q$-$Q$ 图并不是特别能够提供信息，除非样本容量为中等或较大——例如，$n\geqslant20$. 对于小样本，即使在已知观测值是来自正态总体时，$Q$-$Q$ 图的直线性也可能存在很大变异.

\noindent\textbf{例 4.10}\quad (放射性数据的 $Q$-$Q$ 图)

一个微波炉制造厂家的质量控制部门根据联邦政府的要求，监控当炉门关闭着时的辐射量. 测量了随机选取的 $n=42$ 个微波炉在炉门关闭时的辐射量. 数据列于\hyperref[tab:4.1]{表 4.1} 中.

\begin{mvtable}
\phantomsection\label{tab:4.1}\mvtabcap{表 4.1\quad 辐射量数据（关门）}
\mvfit{%
\begin{tabular}{|c|c|c|c|c|c|}
\hline
炉号 & 辐射量 & 炉号 & 辐射量 & 炉号 & 辐射量\\
\hline
1 & 0.15 & 15 & 0.10 & 29 & 0.08\\
2 & 0.09 & 16 & 0.10 & 30 & 0.18\\
3 & 0.18 & 17 & 0.02 & 31 & 0.10\\
4 & 0.10 & 18 & 0.10 & 32 & 0.20\\
5 & 0.05 & 19 & 0.01 & 33 & 0.11\\
6 & 0.12 & 20 & 0.40 & 34 & 0.30\\
7 & 0.08 & 21 & 0.10 & 35 & 0.02\\
8 & 0.05 & 22 & 0.05 & 36 & 0.20\\
9 & 0.08 & 23 & 0.03 & 37 & 0.20\\
10 & 0.10 & 24 & 0.05 & 38 & 0.30\\
11 & 0.07 & 25 & 0.15 & 39 & 0.30\\
12 & 0.02 & 26 & 0.10 & 40 & 0.40\\
13 & 0.01 & 27 & 0.15 & 41 & 0.30\\
14 & 0.10 & 28 & 0.09 & 42 & 0.05\\
\hline
\end{tabular}
}
\end{mvtable}
\noindent\rule{\linewidth}{0.4pt}\par
{\footnotesize 资料来源：数据经 J. D. Cryer 许可使用.}
